\section{Surfaces including m-chains}

\begin{definition}
    An $m$-chain is a union of $m$ lines $l_1, l_2, \dots, l_m$ such that $l_i$ intersects $l_{i+1}$ for $1 \leq i < m$. We will write $C_m = \bigcup_{i=1}^m l_i$ and denote the intersection points as $p_i = l_i \cap l_{i+1}$ for $1 \leq i < m$. A \emph{regular} $m$-chain is an $m$-chain such that no line $l_i$ intersects any $l_j$ for $|i - j| > 1$, and $l_i \neq l_j$ for $i \neq j$.
\end{definition}
\todo[inline]{Improve regular definition wording.}

Notice that an $m$-chain is an $(m+1)$-gon excluding the line $l_m$ intersecting $l_1$. One can then modify the code above ("as in" cite appendix)(!?) to create $m$-chains instead of $(m+1)$-gons, by excluding the final pair of points. The results from this are similar to those of $m$-gons, in that for higher degrees, we find smooth surfaces, up until we find no surfaces. There are however interesting exceptions for degree 2 and 3 surfaces. Here we find singular surfaces when considering 4-chains and 6-chains respectively. We will handle the cases separately. First, we will study the set of $m$-chains in $\mathbb{P}^3$. We will mostly use this in the case of 6-chains, but the results hold in general. (Do we want to use this for 4-chains also?)(!?)
\todo[inline]{Include m2 results for $m$-chains}

\begin{definition}
    We define $Q_m \subseteq \mathbb{G}(1,3)^m$ to be the subset of m-tuples of lines $(l_1, \dots, l_m)$ such that they form an $m$-chain.
\end{definition}

\begin{lemma} \label{lem:dimQm}
    $Q_m$ is an irreducible variety of dimension $3m + 1$ for $m \geq 1$ (Do we have to require $m \geq 1$, or is it obvious!?). 
\end{lemma} 

\todo[inline]{Irreducibility is discussed in EH 3.3.1, saying Schubert cells and their closures are irreducible.}

% \begin{proof}
%     For $m=1$, we have that $Q_1 = \mathbb{G}(1,3)$, which is irreducible of dimension 4. Now, assume that the statement holds for $m - 1$. Consider the projection map $\pi: Q_m \to Q_{m-1}$ given by $\pi(l_1, \dots, l_m) = (l_1, \dots, l_{m-1})$. The fiber over a point $(l_1, \dots, l_{m-1}) \in Q_{m-1}$ is given by the set of lines $l_m$ such that $l_m$ intersects $l_{m-1}$. The set of lines intersecting a given line in $\mathbb{P}^3$ is a 2-dimensional variety (it can be identified with the projective plane of lines through a point in the plane orthogonal to the given line). Thus, the fiber has dimension 2. By the induction hypothesis, we have that $\dim Q_{m-1} = 3(m-1) + 1 = 3m - 2$. Therefore, by the fiber dimension theorem, we have that
%     \[
%     \dim Q_m = \dim Q_{m-1} + \dim \pi^{-1}(l_1, \dots, l_{m-1}) = (3m - 2) + 2 = 3m.
%     \]
%     To show irreducibility, note that $Q_m$ can be constructed as a sequence of projective bundles over $Q_{m-1}$, each time adding a line that intersects the previous one. Since projective bundles over irreducible varieties are irreducible, it follows by induction that $Q_m$ is irreducible. Thus, the lemma holds for all $m$.
% \end{proof}

\begin{proof}
    We observe that $Q_m$ can be constructed as a product of Grassmannians, with a condition of intersection added for $m \geq 2$. That is, if we are given a line $L \in \mathbb{G}(1,3)$, we can consider the space of lines intersecting $L$, denoted as $\Sigma_L = \{l \in \mathbb{G}(1,3) \mid l \cap L \neq \emptyset\}$. Then $Q_m = \mathbb{G}(1,3) \times \Sigma_L^{m-1}$. For m = 1, we have that $Q_1 = \mathbb{G}(1,3)$, which is irreducible of dimension 4. Now, assume that the statement holds for $m-1$. Consider the projection map $\pi: Q_m \to Q_{m-1}$ given by $\pi(l_1, \dots, l_m) = (l_1, \dots, l_{m-1})$. The fiber over a point is given by $\pi^{-1}(l_1, \dots, l_{m-1}) = \{l \in \mathbb{G}(1,3) \mid l \cap l_{m-1} \neq \emptyset\}$. 

    \todo[inline]{Should we denote the fiber over a point as the whole m-tuple, or just the last line? I think the first is more correct.}

    Note that for a given line $L$, $\Sigma_L$ can be constructed as the image of the incidence correspondence $\Gamma = \{(l, p)\in \mathbb{G}(1,3) \times L \mid p \in l\}$ to $\mathbb{G}(1,3)$. The fiber of the projection map $f: \Gamma \to L$ is given by the set of lines through a point $p \in L$, which is isomorphic to $\mathbb{P}^2$. Then the fiber of each point in $L$ is irreducible of dimension 2, and since $f$ is proper, we have that $\Gamma$ is irreducible of dimension 3, meaning that $\Sigma_L$ is irreducible of dimension 3. Thus, $Q_m$ is the product of $m-1$ irreducible varieties of dimension 3, and one Grassmannian of dimension 4, so $Q_m$ is irreducible of dimension $4 + 3(m-1) = 3m + 1$. This completes the proof.
\end{proof}

\todo[inline]{Update proof with notes from 09.10. Don't think this is rigorous enough to show irreducibility.}

\subsection{The case of 4-chains}

\begin{proposition}
    For a general 4-chain in $\mathbb{P}^3$, we cannot find a smooth quadric surface $ S_2 \subset \mathbb{P}^3$ such that it includes the 4-chain.
\end{proposition}

\begin{proof}
    First, assume that $S_2$ is smooth. Since $S_2$ is a smooth quadric surface in $\mathbb{P}^3$, it is of the form $S_2 = \mathbb{P}^1 \times \mathbb{P}^1$. \todo{Cite Beauville} This means that $S_2$ has two rulings, meaning that there are two families of lines on $S_2$, such that any two lines in the same family do not intersect, and any two lines in different families intersect at exactly one point. Denote the two families as $F_1$ and $F_2$. Then we can see that a 4-chain must consist of two lines from $F_1$ and two lines from $F_2$, since any two lines in the same family do not intersect. This means that the 4-chain must be of the form $l_1, l_2, l_3, l_4$ where $l_1, l_3 \in F_1$ and $l_2, l_4 \in F_2$. However, this means that $l_4$ intersects $l_1$, since they are in different families, meaning that the 4-chain is in fact a 4-gon. This is a contradiction, since this is not a general 4-chain. Thus, we cannot find a smooth quadric surface $ S_2 \subset \mathbb{P}^3$ such that it includes the 4-chain.
\end{proof}

\begin{proposition}
    For a general 4-chain in $\mathbb{P}^3$, we can find a singular quadric surface $ S_2 \subset \mathbb{P}^3$ such that it includes the 4-chain.
\end{proposition}
\begin{proof}
    We pair the lines of the 4-chain as $(l_1, l_2)$ and $(l_3, l_4)$. Each pair of lines defines a plane in $\mathbb{P}^3$, denoted as $H_1 = \mathrm{span}(l_1, l_2)$ and $H_2 = \mathrm{span}(l_3, l_4)$. Since the 4-chain is general, we have that $H_1 \neq H_2$. Then the union of the two planes $S_2 = H_1 \cup H_2$ is a singular quadric surface in $\mathbb{P}^3$ that includes the 4-chain.
\end{proof}
\todo[inline]{Improve, used GPT here.}

\subsection{The case of 7-chains}

Now we will look at the case of 7-chains in cubic surfaces.

\begin{proposition}
    For a general 7-chain in $\mathbb{P}^3$, we cannot find a smooth cubic surface $ S_3 \subset \mathbb{P}^3$ such that it includes the 7-chain.
\end{proposition}

% \begin{proof}
%     First, the dimension of the space of cubic surfaces in $\mathbb{P}^3$ is given by the number of monomials of degree 3 in 4 variables, which is $\binom{3+4-1}{4-1} - 1 = \binom{6}{3} - 1 = 19$. Thus, the space of cubic surfaces in $\mathbb{P}^3$ has dimension 19. 

%     Now, the linear system of 7-chains in $\mathbb{P}^3$ has dimension decided by the number of conditions imposed by the 7 lines. Recall that the set of lines in $\mathbb{P}^3$ is parametrized by the Grassmannian variety $\mathrm{Gr}(2,4)$, which is a 4-dimensional variety. Therefore the first line imposes 4 conditions, the second line imposes 3 conditions (since it must intersect the first line), the third line imposes 3 conditions (since it must intersect the second line), and so on. Thus, the total number of conditions imposed by the 7 lines is $4 + 3 + 3 + 3 + 3 + 3 + 3 = 22$. This being of a higher dimension than the space of cubic surfaces, we do not expect to find a cubic surface including our 7-chain.
% \end{proof}

\begin{proof}
    Let $\mathcal{M}$ be the space of smooth cubic surfaces in $\mathbb{P}^3$. This is an open subset of $\mathbb{P}(\mathbb{C}[x_0,\dots, x_3]_3) \cong \mathbb{P}^{19}$. \todo{Why open?} Since $\mathcal{M}$ is an open subset, it shares the dimension, i.e. $\dim \mathcal{M} = 19$. Let $ \mathcal{N} = \{(S_d, l_1, \dots, l_7) \in \mathcal{M} \times \mathbb{G}(1,3)^7 \mid \mathcal{L} \subset S_d, \mathcal{L} \text{ is an $m$-chain}\}$ be the incidence variety of smooth cubic surfaces including a general 7-chain. We have the following diagram:

    \begin{equation}
        \begin{tikzcd}
            & \mathcal{N} \arrow[dl, "\pi_1"'] \arrow[dr, "\pi_2"] & \\
            \mathcal{M}  & & \mathbb{G}(1,3)^7
        \end{tikzcd}. 
    \end{equation}
    \todo[inline]{Should we use $\mathcal{Q}_7$ instead of $\mathbb{G}(1,3)^7$?}
    
    The projection map $\pi_1: \mathcal{N} \to \mathcal{M}$ is given by $\pi_1(S_d, l_1, \dots, l_7) = S_d$. Then the fiber over a point $S_d \in \mathcal{M}$ is given by the set of 7-chains included in $S_d$. Since a smooth cubic surface includes exactly 27 lines,\todo{Refer to result} we can find a 7-chain by choosing 7 lines from these 27. We know from \cref{prop:12-gon}(!?) \todo{Result in this section} that it is possible to find a 7-chain, and so this fiber is not empty. Furthermore, there are only a finite number of 7-chains when there is a finite number of lines. Then the dimension of the fiber $\dim \pi_1^{-1}(S_d)$ is 0, being a finite set of 7-chains. Thus, from \cref{thm:fiber}, we have that $\dim \mathcal{N} = \dim \mathcal{M} + \dim \pi_1^{-1}(S_d) = 19 + 0 = 19$. \todo[inline]{Refer to theorem about dimension of fiber.}

    Now we are interested in the fiber of the set of 7-chains in $\mathbb{P}^3$, $\mathcal{Q}_7 \subseteq \mathbb{G}(1,3)^7$. By \cref{lem:dimQm}, we have that $\dim \mathcal{Q}_7 = 22$, and so the projection map $\pi_2: \mathcal{N} \to \mathbb{G}(1,3)^7$ given by $\pi_2(S_d, l_1, \dots, l_7) = (l_1, \dots, l_7)$ cannot be surjective on (dominate!?) $\mathcal{Q}_7$. This means that taking a general 7-chain $C_7 \in \mathcal{Q}_7$, and considering the fiber $\pi_2^{-1}(C_7)$, we do not expect this to be non-empty, meaning we will not have a surface including it. Thus, for a general 7-chain in $\mathbb{P}^3$, we cannot find a smooth cubic surface $S_3 \subset \mathbb{P}^3$ such that it includes the 7-chain.
\end{proof}

\todo[inline]{Proof should include: projection is surjective (every smooth cubic surface has at least a 7-chain, shown in this section), implies dim N $\geq$ dim M. Then fiber is finite, so dim N $\leq$ dim M. Prop 10.26?}

Again, as in the case for 4-chains in quadric surfaces, computations show that we can find a singular cubic surface including a general 7-chain.

\begin{proposition}
    For a general 7-chain in $\mathbb{P}^3$, we can find a \emph{singular} cubic surface $ S_3 \subset \mathbb{P}^3$ such that it includes the 7-chain. (!?)
\end{proposition}
\begin{proof}
    Find a nice example of a singular cubic surface including a general 7-chain.
\end{proof}

\todo[inline]{Need to rework entire section above, should be 6-chains, not 7-chains.}

Note that if we instead asked whether a smooth cubic surface in $\mathbb{P}^3$ can include a 7-chain, the situation is different. In this case, we can indeed find a 7-chain such that it is included in the surface. In fact, we know we could find an 11-chain in this case, and a 12-gon. We show this in the following proposition. 

\todo[inline]{Maybe move this proposition earlier, to use in proof. Maybe lemma? Also interesting in itself, so maybe not.}
\begin{proposition}\label{prop:12-gon}
    For a smooth cubic surface in $\mathbb{P}^3$, we can find a 12-gon such that it is included in the surface.
\end{proposition}

\begin{proof}
    Recall from \cref{prop:blowup6} that a smooth cubic surface in $\mathbb{P}^3$ is isomorphic to the blow-up of $\mathbb{P}^2$ in 6 points. It may be useful to again consider the picture in figure \cref{fig:blowup6} on page (!?). The surface includes 27 lines, which are precisely the exceptional divisors of the blow-up, the strict transforms of the 15 lines through pairs of points, and the strict transforms of the 6 conics through 5 of the points (\cref{prop:27}). We can then find a 12-gon by considering the following lines: Let $E_1, E_2, E_3, E_4, E_5, E_6$ be the exceptional divisors, let $L_{ij}$ be the strict transform of the line through points $p_i$ and $p_j$, and let $C_i$ be the strict transform of the conic through all points except $p_i$. Then we can consider the 12-gon given by the lines $E_1, L_{12}, E_2, L_{23}, E_3, L_{34}, E_4, L_{45}, E_5, L_{56}, E_6, C_1, L_{16}$. This is indeed a (general!?) 12-gon since each line intersects exactly two other lines in the set.
\end{proof}

By removing one of the lines from the 12-gon, we immediately get the following:
\begin{corollary}
    For a smooth cubic surface in $\mathbb{P}^3$, we can find an 11-chain such that it is included in the surface.
\end{corollary}

This implies that the $m$-chains longer than 7 included in the smooth cubic surfaces are somehow special. \todo{Fill out later.}

\subsection{The case of 6-chains}

Now we will examine 6-chains in cubic surfaces. We will show that a general 6-chain is not included in a smooth cubic surface, but that we can find singular cubic surfaces including a general 6-chain. We start by studying the incidence variety of smooth cubic surfaces including a general 6-chain.

Let $\mathcal{M}$ be the set of smooth cubic surfaces in $\mathbb{P}^3$. This is an open subset of $\mathbb{P}(\mathbb{C}[x_0,\dots, x_3]_3) \cong \mathbb{P}^{19}$. \todo{Why open?} Since $\mathcal{M}$ is an open subset, it shares the dimension, i.e. $\dim \mathcal{M} = 19$. Let $ \mathcal{N} = \{(S_d, l_1, \dots, l_6) \in \mathcal{M} \times \mathbb{G}(1,3)^6 \mid \mathcal{L} \subset S_d, \mathcal{L} \text{ is an $m$-chain}\}$ be the incidence variety of smooth cubic surfaces including a general 6-chain. We have the following diagram:

\begin{equation}
    \begin{tikzcd}
        & \mathcal{N} \arrow[dl, "\pi_1"'] \arrow[dr, "\pi_2"] & \\
        \mathcal{M}  & & \mathbb{G}(1,3)^6
    \end{tikzcd}. 
\end{equation}
\todo[inline]{Should we use $\mathcal{Q}_7$ instead of $\mathbb{G}(1,3)^7$?}

\begin{remark}
    Note that it is important that we consider the incidence variety of over specifically the smooth cubic surfaces, including a general 6-chain. If not, one can notice that the fiber of the singular surfaces which are the union of 3 planes, will include a much higher dimensional family of 6-chains, and so the incidence variety will not be irreducible. (!?)
\end{remark}

\begin{lemma}
    For a smooth cubic surface, we cannot find a regular(!?) 6-chain. 
\end{lemma}
\begin{proof}
    We denote the 27 lines of the cubic surface as $C_1, \dots, C_27$. Let $Q_5$ be a regular 5-chain, consisting of $C_1, \dots, C_5$. Look at the class $D = C_1 + C_2 - C_4 - C_5$. Then $D^2 = 0$, and $HD=0$. The Hodge index theorem(!?) tells us that $D=0$, and we have the equality $C_1 + C_2 = C_4 + C_5$. Now we know that if a line $C_i$ for $i \neq 4$ (then $C_4^2 + C_4C_5 = 0$) intersects $C_5$, it must also intersect either $C_1$ or $C_2$, and would not make a regular 6-chain. 
\end{proof}